\paragraph{Sharpness-aware minimization.}
Flat minimizers are more robust to the shifts between training and test losses, thereby enhancing the generalization ability of neural networks ~\citep{DBLP:conf/iclr/JiangNMKB20, DBLP:conf/nips/PetzkaKASB21, DBLP:conf/uai/DziugaiteR17}. The relationship between generalization and the width of minima has been studied both theoretically and empirically in several prior works ~\citep{DBLP:conf/nips/HochreiterS94, neyshabur2017exploring, dinh2017sharp, fort2019emergent}. Consequently, a variety of methods have been developed to search for flat minima ~\citep{DBLP:conf/iclr/PereyraTCKH17, Chaudhari2017EntropySGDBG, keskar2017large, DBLP:conf/uai/IzmailovPGVW18}. Among the flat minimizers, Sharpness-Aware Minimization (SAM) \cite{sam} has gained significant attention due to its effectiveness. SAM has been leveraged across a wide range of tasks and domains, including domain generalization~\citep{cha2021swad, wang2023sharpness, zhang2023flatness}, federated learning~\citep{caldarola2022improving, qu2022generalized}, Bayesian networks~\citep{va2023_flat_Bayes, llenhoff2023sam}, and meta-learning~\citep{abbas2022sharp}. Moreover, SAM has demonstrated its ability to enhance generalization in both vision models~\citep{chen2021vision} and language models~\citep{bahri-etal-2022-sharpness}.

Nevertheless, these studies are constrained to finite-dimensional Euclidean spaces. In this work, we strengthen these generalization principles to infinite-dimensional functional spaces and propose a particle-sampling method grounded in this theoretical framework.

\paragraph{Bayesian Inference.} Two main strategies were widely employed in the literature of Bayesian inference. The first paradigm is \textit{Variational Inference}, which aims to approximate a target distribution by selecting a distribution from a family of potential approximations and optimizing a variational lower bound. \citet{graves2011practical} introduced the use of a Gaussian variational posterior approximation for neural network weights, which was later extended by \citet{kingma2013auto, kingma2015variational, blundell2015weight} with the reparameterization trick to facilitate training deep latent variable models. \citet{louizos2017multiplicative} proposed using a matrix-variate Gaussian to model entire weight matrices \citep{gupta2018matrix} to increase further the flexibility of posterior approximations, which offers a novel approach to approximate the posterior. Subsequently, various alternative structured forms of the variational Gaussian posterior were proposed, including the Kronecker-factored approximations \citep{zhang2018noisy, ritter2018scalable, rossi2020walsh}, or non-centered or rank-1 parameterizations \citep{ghosh2018structured, dusenberry2020efficient}. Recently, \citet{quan2025promoting} proposes to incorporate sharpness and distributional robustness to Bayesian Inference, thereby modelling the interactions between particles and allowing a more diverse set of model particles. 

The second paradigm of Bayesian inference is \textit{Markov Chain Monte Carlo} (MCMC), which involves sampling multiple models from the posterior distribution. MCMC has been applied to neural network inference, such as Hamiltonian Monte Carlo (HMC) \citep{Neal_Bayesian}. However, HMC requires the computation of full gradients, which can be computationally expensive. To address this, Stochastic Gradient Langevin Dynamics (SGLD) \citep{sgld} integrates first-order Langevin dynamics within a stochastic gradient framework. Stochastic Gradient HMC (SGHMC) \citep{chen2014stochastic} further incorporates stochastic gradients into Bayesian inference, enabling scalability and efficient exploration of different solutions. Another critical approach, Stein Variational Gradient Descent (SVGD) \citep{svgd}, closely related to our work, uses a set of particles that converge to the target distribution. It is also theoretically established that SGHMC, SGLD, and SVGD asymptotically sample from the posterior as the step sizes approach zero.